% EDITING: type inside the final {...} of each field. Change the shown font size
% (for example, font=10pt) in that same field block when needed.

\GrantSection{4}{Clinical Trial Synopsis}

\GrantGuidance{Complete this module when a clinical trial is proposed or when a funder requests a protocol synopsis. It is separate from the five-page Pioneer-style essay so that it can be included, omitted, or converted into the target funder's study-record format.}

% -----------------------------------------------------------------------------
% Page-local wrapped table field.
%
% This definition applies only inside the following group and does not modify
% clinicaltrialgrant.sty or any other page in the document.
% -----------------------------------------------------------------------------
\begingroup
\ExplSyntaxOn
\NewDocumentCommand{\GrantWrappedTableField}{ O{} +m }
  {
    \grant_field_setup:nnnnn
      {#1}
      {0.98\linewidth}
      {0.48in}
      {120}
      {
        \TextField[
          name={\GrantResolvedID},
          value={#2},
          multiline=true,
          width=\GrantResolvedWidth,
          height=\GrantResolvedHeight,
          charsize=\GrantResolvedFont,
          maxlen=\GrantResolvedMaxChars,
          bordercolor={\GrantPDFBorderColor},
          backgroundcolor={\GrantPDFBackgroundColor},
          borderwidth=1
        ]{}
      }
  }
\ExplSyntaxOff

\begin{tabularx}{\linewidth}{
  >{\bfseries\raggedright\arraybackslash}p{0.29\linewidth}
  >{\raggedright\arraybackslash}X
}

Protocol title &
\GrantWrappedTableField[
  id=trial-protocol-title,
  font=9pt,
  width=0.98\linewidth,
  height=0.5in,
  maxchars=220
]{On-Premises LLM-Directed Robotic Pancreaticoduodenectomy with Perioperative Daraxonrasib (RMC-6236) in KRAS-Mutated Pancreatic ^^JDuctal Adenocarcinoma \citep{onpremwhippl}
}\\[5pt]

Protocol / version / date &
\GrantTableField[
  id=trial-protocol-number,
  font=9pt,
  width=0.98\linewidth,
  maxchars=120 
]{Protocol v1.0, Repository v4.2.0, June 29, 2026.}\\[5pt]

Phase &
\GrantTableField[
  id=trial-phase,
  font=9pt,
  width=0.98\linewidth,
  maxchars=40
]{Phase 1 \citep{phase1ind,onpremwhippl}
}\\[5pt]

Study type &
\GrantWrappedTableField[
  id=trial-type,
  font=9pt,
  width=0.98\linewidth,
  height=0.36in,
  maxchars=100
]{First-in-human, prospective, open-label, single-arm interventional clinical study \citep{phase1ind}
}\\[5pt]

Condition / indication &
\GrantWrappedTableField[
  id=trial-condition,
  font=9pt,
  width=0.98\linewidth,
  height=0.35in,
  maxchars=160
]{Resectable or borderline-resectable KRAS G12-mutated pancreatic ductal adenocarcinoma (PDAC) \citep{phase1ind}
}\\[5pt]

Investigational intervention &
\GrantWrappedTableField[
  id=trial-intervention,
  font=9pt,
  width=0.98\linewidth,
  height=0.4in,
  maxchars=180
]{Perioperative daraxonrasib (RMC-6236) administered with an LLM-directed eight-arm robotic pancreaticoduodenectomy (Whipple) \citep{phase1ind,onpremwhippl}
}\\[5pt]

Comparator / control &
\GrantWrappedTableField[
  id=trial-comparator,
  font=9pt,
  width=0.98\linewidth,
  height=0.17in,
  maxchars=160
]{None for the Phase 1 single-arm study; comparative evaluation is deferred to the later randomized Phase 2 study \citep{phase1ind,daraxwhipple}
}\\[5pt]

Allocation / masking &
\GrantTableField[
  id=trial-allocation,
  font=9pt,
  width=0.98\linewidth,
  height=0.17in,
  maxchars=120
]{Phase I: Non-randomized, Open-label, Single-arm.}\\[5pt]

Target enrollment &
\GrantTableField[
  id=trial-enrollment,
  font=9pt,
  width=0.98\linewidth,
  height=0.17in,
  maxchars=40
]{Phase I: 13 patients}\\[5pt]

Number and type of sites &
\GrantWrappedTableField[
  id=trial-sites,
  font=9pt,
  width=0.98\linewidth,
  height=0.15in,
  maxchars=140
]{Local study: Single-center; NCI-designated Comprehensive Cancer Center.}\\[5pt]

Study duration / follow-up &
\GrantWrappedTableField[
  id=trial-duration,
  font=9pt,
  width=0.98\linewidth,
  height=0.25in,
  maxchars=100
]{Short-term: 24 months total (12-month accrual, 12-month follow-up).}

\end{tabularx}

\endgroup

\GrantTextArea[
  id=trial-rationale,
  font=10pt,
  height=0.9in,
  maxchars=1400
]{Scientific premise and clinical rationale}{PDAC has poor long-term survival, and even intended curative resections can be limited by perioperative morbidity or failure to achieve margin-negative surgery. The public Phase 1 and Phase 2 records frame the intervention as perioperative daraxonrasib combined with an LLM-directed robotic Whipple in KRAS-mutated PDAC, with the later Phase 2 protocol explicitly posing the comparative efficacy question after initial feasibility and dose-setting work. \citep{phase1ind,onpremwhippl,daraxwhipple}
}

\GrantTextArea[
  id=trial-primary-endpoint,
  font=10pt,
  height=2.25in,
  maxchars=1300
]{Primary objective, endpoint, estimand, and time point}{Primary objective: to estimate the safety, tolerability, and operational feasibility of perioperative daraxonrasib administered with an LLM-directed robotic Whipple procedure in adults with resectable or borderline-resectable KRAS G12-mutated PDAC, and to identify a regimen and operating context ^^Jsuitable for later-phase testing \citep{phase1ind,onpremwhippl}. Primary endpoint: the participant level incidence of protocol-defined dose-limiting toxicity, treatment-related serious adverse events, unplanned major procedure-related morbidity, or predefined cyber-physical safety failure within ^^Jthe primary assessment window. Primary estimand: the proportion of enrolled participants who, under ^^Jthe assigned perioperative daraxonrasib-plus-robotic strategy, experience the composite primary endpoint regardless of temporary interruption, dose reduction, conversion to open surgery, or rescue actions; deaths and clinically significant safety events occurring after discontinuation but within the assessment window will be counted. Primary time point: first study treatment through 30 days after index surgery, with 90-day postoperative events summarized as supportive safety context.
}

\GrantTextArea[
  id=trial-secondary-endpoints,
  font=10pt,
  height=2.32in,
  maxchars=1250
]{Key secondary and exploratory objectives / endpoints}{Key secondary objectives are to characterize the safety profile by attribution, grade, seriousness, and timing; to estimate feasibility of completing the planned robotic procedure under the governed operating framework; to describe perioperative recovery and 90-day surgical outcomes; and to assess whether the regimen supports timely transition to standard postoperative care \citep{phase1ind,onpremwhippl,daraxwhipple}. Key secondary endpoints therefore include rates of treatment-emergent adverse ^^Jevents, serious adverse events, Clavien-Dindo grade III or higher complications, conversion to open surgery, inability to complete the planned robotic workflow, hospital length of stay, readmission, reoperation, margin status, and time to adjuvant-therapy readiness. Exploratory objectives are to ^^Jdescribe early signals in ctDNA or other protocol-specified biomarkers, radiographic tumor response or stability, pathologic treatment effect, pharmacodynamic readouts, robotic performance logs, verification gate disposition rates, and documentation-efficiency metrics from the repository-based trial workflow \citep{phase1guidance,oncotrialllm,qspmetpancre,pdacdigtwinp,vvuqllmverif}.
}

\GrantTextArea[
  id=trial-eligibility,
  font=10pt,
  height=2.28in,
  maxchars=1400
]{Eligibility criteria and target population}{Target population: adults with resectable or borderline-resectable pancreatic ductal adenocarcinoma harboring a KRAS G12 alteration for whom curative-intent pancreaticoduodenectomy is clinically appropriate \citep{phase1ind,onpremwhippl}. Key inclusion criteria are age 18 years or older; ^^Jhistologic or cytologic evidence of PDAC or unequivocal clinicopathologic diagnosis per protocol; documented KRAS G12-mutated status from a CLIA-qualified or equivalently validated assay; ECOG performance status 0 to 1; adequate marrow, hepatic, renal, and cardiopulmonary function; resectable or borderline-resectable disease without distant metastasis; anatomy suitable for the planned robotic approach; and ability to understand and sign informed consent. Key exclusion criteria are metastatic disease, uncontrolled infection, clinically significant intercurrent illness, prior intolerance or ^^Jcontraindication to study drug or major pancreatic surgery, pregnancy or breastfeeding, inability to comply with perioperative monitoring, and any medical, technical, or psychosocial condition that, in the investigator's judgment, would make study participation unsafe or uninterpretable.
}

